\section{Regular receipt lifts}

Let \(X=(V,E)\) be a connected finite graph and let \(R\) be a finite
group.  Choose an orientation of each undirected edge and assign a
voltage
\[
  \alpha(u,v)\in R,
  \qquad
  \alpha(v,u)=\alpha(u,v)^{-1}.
\]
The associated regular lift has vertex set \(V\times R\), with
\[
  (u,r)\sim (v,r\alpha(u,v)).
\]
Changing the section of the fibers applies a vertex gauge.  For
\(g:V\to R\), the transformed voltage is
\[
  \alpha^g(u,v)
  =
  g(u)^{-1}\alpha(u,v)g(v).
\]
Thus individual edge values are chart-dependent, while circuit
holonomy is gauge-covariant.

Fix a root \(x_0\) and a spanning tree \(T\).  Every voltage assignment
is gauge-equivalent to one that is trivial on \(T\).  The remaining
chord voltages give the fundamental circuit receipts.  They determine
a homomorphism
\[
  \rho_\alpha:\pi_1(X,x_0)\longrightarrow R
\]
up to conjugation in \(R\).

\begin{theorem}[Finite receipt classification]
Gauge classes of regular \(R\)-receipt lifts of a connected graph
\(X\) correspond to \(R\)-conjugacy classes of homomorphisms
\(\pi_1(X)\to R\).  If \(H\) is the image, then the lift has
\([R:H]\) connected components, each containing
\(\lvert V(X)\rvert\lvert H\rvert\) vertices.  In particular, the lift
is connected exactly when \(H=R\).
\end{theorem}

For \(R=\Ctwo\), a nontrivial circuit receipt doubles the visible return
cycle.  For \(R=\Vfour\), connectedness requires two independent
binary holonomy directions.  These statements classify a proposed
lift; they do not select the proposal.  The native encounter begins at
that boundary.
